\section{Устройство робота}

Работа состоит из двух программ. Обе принимают единственный аргумент --- путь до
общего yaml-конфига.

\begin{itemize}
    \item \texttt{crawler.py} --- обкачка: берёт URL из очереди, скачивает
          документ, выделяет текст, сохраняет в базу и пополняет очередь
          найденными ссылками;
    \item \texttt{recrawler.py} --- переобкачка: проверяет, не изменились ли уже
          скачанные документы, и изменившиеся ставит в очередь.
\end{itemize}

Общий код вынесен в модуль \texttt{common.py}: конфиг, нормализация URL,
определение кодировки, выделение текста, хранилище и HTTP-клиент.

\subsection*{Хранилище}

В качестве базы данных выбран SQLite. Задание называет MongoDB как пример, но
SQLite решает ту же задачу и не требует поднимать сервер: модуль
\texttt{sqlite3} входит в стандартную библиотеку, а вся база --- один файл,
который удобно приложить к отчёту.

В базе две таблицы. \texttt{documents} хранит сами документы:

\begin{itemize}
    \item \texttt{url} --- нормализованный адрес, уникальный ключ;
    \item \texttt{raw\_html} --- \enquote{сырой} html-текст документа;
    \item \texttt{text} --- выделенный текст без навигационной обвязки;
    \item \texttt{source} --- название источника;
    \item \texttt{crawled\_at} --- дата обкачки, Unix timestamp;
    \item \texttt{content\_hash} --- отпечаток текста, нужен переобкачке;
    \item \texttt{etag}, \texttt{last\_modified} --- для условных запросов;
    \item \texttt{revision} --- сколько раз документ переобкачивался.
\end{itemize}

Задание требует хранить \enquote{сырой} HTML, а лекция --- только необходимый
текст. Требования не противоречат друг другу, если хранить и то, и другое:
\texttt{raw\_html} пригодится, если правила выделения текста изменятся, а
переразбирать корпус по ним, не выкачивая его заново, гораздо дешевле.

Вторая таблица, \texttt{frontier}, --- это очередь URL: адрес, источник, глубина,
статус (\texttt{pending}, \texttt{done}, \texttt{failed}, \texttt{skipped}) и
приоритет.

\subsection*{Остановка и продолжение}

Очередь лежит в базе, а не в памяти. URL помечается как \texttt{done} только
после того, как документ сохранён, поэтому прерывание в любой момент безопасно:
недообработанный адрес остаётся в статусе \texttt{pending} и при следующем
запуске будет взят снова. Сигналы \texttt{SIGINT} и \texttt{SIGTERM}
перехватываются --- робот дописывает текущий документ и выходит, а не умирает
посреди записи.

Уникальность \texttt{url} в обеих таблицах снимает вопрос дублирования ссылок:
повторное добавление просто игнорируется базой.

\subsection*{Нормализация URL}

Часть ссылок на странице относительные, поэтому при разборе они склеиваются с
адресом текущей страницы. Дальше из адреса убирается всё, что не влияет на
содержимое: якорь, служебные параметры (\texttt{utm\_*}, \texttt{gclid} и
подобные), порт по умолчанию, регистр схемы и хоста; параметры запроса
сортируются, повторные слэши схлопываются.

Отдельно пришлось раскрывать процентные последовательности. Ссылки в HTML
Википедии закодированы, и адрес \texttt{/wiki/\%D0\%A4\%D0\%B8\%D0\%BB\%D1\%8C\%D0\%BC}
--- это тот же документ, что и \texttt{/wiki/Фильм}. Без раскрытия в базе
появлялись бы два документа вместо одного.

\subsection*{Определение кодировки}

Порядок ровно такой, как на лекции: сначала заголовок \texttt{Content-Type},
затем \texttt{<meta charset>} в самом HTML, и только потом UTF-8 по умолчанию.
Это не теоретическое требование: один из двух источников,
\texttt{kino-teatr.ru}, отдаёт страницы в windows-1251, и без определения
кодировки весь текст с него попал бы в базу нечитаемым.

\subsection*{Вежливость}

Робот соблюдает две вещи. Явные правила --- \texttt{robots.txt}: он скачивается
для каждого хоста и кэшируется, запрещённые адреса помечаются как
\texttt{skipped}. Неявные --- пауза между запросами к одному хосту, её величина
задаётся в секции \texttt{logic} конфига.

\subsection*{Выбор следующего URL}

Очередь отдаёт документы так: сначала всё, что ждёт переобкачки, затем ---
источник, отличный от предыдущего. Чередование источников нужно по двум
причинам. Пока идёт запрос к одному хосту, для другого успевает истечь пауза
вежливости, и робот не простаивает. И оба источника наполняются равномерно, а не
по очереди целыми сайтами.

\subsection*{Фильтры ссылок}

Фильтров два, и это оказалось важным. Первый, \texttt{allow}/\texttt{deny},
решает, обходить ли ссылку вообще. Второй, \texttt{store\_allow} и
\texttt{store\_deny}, решает, сохранять ли страницу как документ.

Разделять их пришлось потому, что каталоги, списки и страницы категорий нужны
роботу только как источник ссылок. Текста в них хватает, чтобы пройти фильтр по
длине, и в первой версии робота база наполнилась оглавлениями вместо статей.

По той же причине в конфиге есть \texttt{follow\_only} --- список страниц, с
которых вообще собираются ссылки. Без него робот уходит из своей тематики: из
статьи о фильме ссылки ведут на произвольные статьи Википедии, и через несколько
шагов корпус перестаёт быть кинематографическим.

\pagebreak
